\chapter{The mannequin arm and how it commands the robot}
\label{ch:master}

\section{What the master arm is}

The master is a scale model of the robot arm, about a third of the size, with
the same sequence of joints. A rotary sensor sits at each joint and a small
inertial sensor sits at the wrist, and a microcontroller reads them all and
sends the numbers up a USB cable. The driver holds it by the tip and moves it.
Nothing pushes back.

\begin{figure}[htbp]
  \centering
  \includegraphics[width=0.78\linewidth]{master_arm_chain.png}
  \caption[The master arm]{The master arm as modelled, with its joint sequence
  labelled. Its links total \SI{0.272}{\metre} against the robot's
  \SI{0.902}{\metre} of reach, so the same joint angles do not produce the same
  hand movement at each end.}
  \label{fig:master}
\end{figure}

\section{Why the joints are not copied directly}

The first version of this system sent each master joint angle straight to the
matching robot joint. That is the simplest thing to build and it was
abandoned, for a reason that is easy to see once stated: the two arms are not
the same size. A \SI{10}{\degree} bend at the master's elbow moves its tip a
few centimetres and the robot's hand about three times as far. Worse, the
robot's joint limits and the master's are different, so there are master poses
with no robot equivalent at all.

What replaced it maps movement to movement. The driver's hand moving
\SI{100}{\milli\metre} forward asks the robot's hand to move
\SI{100}{\milli\metre} forward, from wherever it happens to be. The robot's
own solver works out the joint angles. This decouples the two arms completely
and costs one thing: the robot's spare joint, the one that swings the elbow
without moving the hand, is left unspecified and has to be chosen by something
else.

\subsection{Lifting the mouse}

Because the master's reach is a third of the robot's, the driver runs out of
arm before the robot runs out of workspace. The fix is the one a computer
mouse uses. A button disengages the link, the driver brings their hand back to
a comfortable position, and pressing again re-engages from there. The robot
does not move during the disengagement.

Getting this right took three separate corrections. The button had to be
debounced, because a single press was registering twice. The reference
position had to be captured only when the master was genuinely still,
otherwise the system latched onto a moving target. And the smoothing filter
had to be cleared on re-engagement, or the first few frames after re-engaging
were a blend of the new position and the old one. Measured over eight
re-engagements, the arm now travels the full distance each time with no
snap-back.

\section{Reading the sensors: what the channels actually do}

The specification says fourteen rotary sensors. What arrives is less tidy.

A failing sensor of this kind does not go quiet. It reports values that span
the full range and are not a trajectory: consecutive readings unrelated to
each other, jumping tens of degrees in a fraction of a second. Judged on range
alone, such a channel looks healthy, which is exactly how a broken one keeps
being trusted. The system therefore classifies each channel into one of four
states, and the third of them had to be invented for this project.

\begin{table}[htbp]
  \centering
  \small
  \caption[Channel verdicts]{The four verdicts, and the current state of the
  master arm. Measured over \num{20430} recorded frames.}
  \label{tab:verdicts}
  \begin{tabular}{@{}p{0.16\linewidth}p{0.50\linewidth}p{0.26\linewidth}@{}}
    \toprule
    \textbf{Verdict} & \textbf{Meaning} & \textbf{Channels} \\
    \midrule
    Working & Reports a smooth range of values, rarely drops out & 5 of 14 \\
    Dropping out & Reports usable values with gaps & 1 of 14 \\
    Incoherent & Present, fresh, and wrong: more than one update in twenty
      jumps by over \SI{60}{\degree} & 6 of 14 \\
    Dead & No usable reading at all & 2 of 14 \\
    \bottomrule
  \end{tabular}
\end{table}

\begin{figure}[htbp]
  \centering
  \includegraphics[width=0.95\linewidth]{channels.pdf}
  \caption[Channel health]{Every channel, on both measures that decide its
  verdict. The upper panel is the one that matters: a channel above the five
  per cent line is reporting values that are present, fresh and not a
  trajectory. Six channels sit above it, five of them on the left arm.
  Produced from the stored channel record of 6 August 2026.}
  \label{fig:channels}
\end{figure}

One correction belongs here because it changed two conclusions that had stood
for weeks. The recorder writes rows at about \SI{50}{\hertz} while the sensors
themselves update at fifteen to seventeen. Between four and five of every six
rows are therefore an exact copy of the row before. Any statistic computed
across consecutive rows is dominated by those copies, which is why the noise
figure read as exactly zero on every channel. It was measuring the recorder.
Computed between genuine sensor updates instead, the noise floor is
\SIrange{0.2}{5.4}{\degree}, and one channel that had been written off as dead
turned out to be the cleanest one on that arm.

\section{Working with fewer channels}

\Cref{fig:ladder} shows what happens as channels are lost. The system does not
have a working state and a broken state. It has a ladder, and it announces
which rung it is standing on.

\begin{figure}[htbp]
  \centering
  % What the system does as sensor channels are lost.
\begin{tikzpicture}[x=1mm, y=1mm, font=\scriptsize]

  \node[sense, minimum height=9mm, text width=38mm] (raw) at (0,0)
    {Fourteen sensor readings arrive, about fifty times a second};

  \node[choice, below=8mm of raw, text width=30mm] (health)
    {Is this channel behaving?};
  \draw[go] (raw) -- (health);

  \node[stage, right=22mm of health, text width=44mm, minimum height=16mm] (verdict)
    {\textbf{One of four verdicts}\\[2pt]
     \emph{working} \textbullet\ \emph{dropping out} \textbullet\
     \emph{incoherent} (present, fresh and wrong) \textbullet\ \emph{dead}};
  \draw[go] (health) -- (verdict);

  \node[stage, below=9mm of health, text width=44mm, minimum height=12mm] (mode)
    {\textbf{Pick the best mapping the surviving channels allow}, and publish
     which one is in use};
  \draw[go] (health) -- (mode);

  \foreach \i/\lbl/\cost in {
      0/{All shoulder channels working}/{full three-dimensional command},
      1/{One shoulder channel lost}/{costs \SI{63.8}{\milli\metre} of tip accuracy},
      2/{Only one channel left}/{direction only, distance held fixed},
      3/{No usable channel}/{\textbf{refuse: report no position}}}
  {
    \pgfmathsetmacro{\y}{-34 - \i*11}
    \ifnum\i=3
      \node[guard, text width=76mm, minimum height=8mm] at (14,\y) {\lbl \ \textbullet\ \cost};
    \else
      \node[stage, text width=76mm, minimum height=8mm] at (14,\y) {\lbl \ \textbullet\ \cost};
    \fi
  }
  \draw[go] (mode) -- (14,-30);

  \node[tag, text width=86mm, align=left] at (14,-80)
    {The refusal at the bottom is the important rung. On a machine mounted
     beside somebody's head, a position that has been guessed is worse than an
     admitted gap, and the difference has to be visible from outside the
     program.};

\end{tikzpicture}

  \caption[The degradation ladder]{What the system does as sensor channels
  fail. Each rung is a different way of estimating where the driver's hand is,
  using fewer inputs than the one above it. The cost of each step down is
  measured rather than assumed, and the bottom rung is a refusal.}
  \label{fig:ladder}
\end{figure}

The mapping used at the top of the ladder splits the driver's hand position
into three parts and takes each from whichever sensor sees it most directly.
How high the hand is comes from the wrist's inertial sensor, which feels
gravity and therefore needs no calibration and no drift correction. Which way
round the hand is comes from the shoulder joint, which is a single rotary
sensor. How far out the hand is comes from the remaining joint angles. The
point of splitting it this way is that losing one part does not destroy the
others.

Two numbers say what this costs. Dropping from four sensed joints to one costs
\SI{63.8}{\milli\metre} of accuracy at the tip, measured against the recorded
frames. And replaying the failing channels with the driver's hand deliberately
held still shows them injecting \SI{54.5}{\milli\metre} of commanded movement
that nobody asked for, against a master reach of \SI{272}{\milli\metre}. The
second figure is the reason the failing channels are frozen rather than
filtered: the arm is a relay, and whatever the sensors feed in, the robot
faithfully reproduces.

\section{Making the motion smooth}

Until late in the project, the last step before the joint commands went out
was a simple limit on how far each joint could move in one cycle. It is the
obvious approach and it has a defect that is not obvious.

Because each joint is limited separately, a joint with a short distance to
travel arrives in one cycle while a joint with a long distance takes two. In
between, the joints are at different fractions of their own journeys, and the
hand is not on the straight line the solver intended. Measured on a typical
move, the hand strayed \SI{51.3}{\milli\metre} from the line it was supposed
to follow. For comparison, the gripper has to be within \SI{30}{\milli\metre}
of an object to close on it.

The fix was to replace the per-joint limit with a generator that plans all
seven joints together, so they start together, finish together, and stay in
proportion in between~\cite{berscheid2021ruckig}. The straying went to zero,
every joint now arrives on the same cycle, and the peak commanded joint speed
dropped from about twelve times the manufacturer's stated limit to exactly
one times it. The old code had never read the limits file at all.

This is worth stating plainly because it was not a tuning problem. Sweeping
the old method's step size across a wide range does not make the straying go
away; it settles at about \SI{68}{\milli\metre} and stays there. The fault was
the independence of the limits, not their size.

\section{What the simulation does not tell us}

Every joint on the real arm stops slightly short of where it was sent, on the
side it approached from. Measured across 252 recorded joint movements, 211 of
them land between \SI{0.25}{\degree} and \SI{0.30}{\degree} short, with the
direction following the direction of travel. That is a bound rather than a
spread, which means it is a systematic offset and not noise.

At the hand, that adds up to about \SI{7.2}{\milli\metre} on a
\SI{100}{\milli\metre} move. A single correction per arm, which simply
overshoots each joint slightly in the direction it is already going, removes
about three quarters of it. The correction is implemented and switched off,
because it has not yet been tried on the physical arm and switching it on
untested would change every recorded figure at once.
